Write \(n=2k\), where \(k\geq2\), and define the orthonormal pair-sum and pair-difference vectors
\[
e_j^+=\frac{e_{2j-1}+e_{2j}}{\sqrt2},
\qquad
e_j^-=\frac{e_{2j-1}-e_{2j}}{\sqrt2},
\qquad j=1,\ldots,k.
\]
Let \(T=\operatorname{span}\{e_1^-,\ldots,e_k^-\}\).  The pair-score definition gives the orthogonal decomposition
\[
H_n=S_{\mathrm{pair}}\mathbin{\oplus}T,
\qquad
\dim(S_{\mathrm{pair}})=k-1,
\qquad
\dim(T)=k.
\tag{1}
\]
In particular, both summands are nonzero.  Let \(G\) be the finite group generated by permutations of the \(k\) pairs and swaps of the two units within any pair, and let \(U_g\) be the permutation matrix of \(g\in G\).  This group preserves \(\mathcal Z_n\), \(H_n\), and \(S_{\mathrm{pair}}\).  Hence, for every \(A\in\mathcal C_n\),
\[
\overline A=\frac1{|G|}\sum_{g\in G}U_gAU_g^\top\in\mathcal C_n.
\tag{2}
\]
The feasible unit-vector set defining \(q_\rho(\,cdot\,;S_{\mathrm{pair}})\) is also invariant under \(G\).  Since \(q_\rho\) is a pointwise maximum of linear functions of its covariance argument, it is convex.  Invariance and Jensen's inequality therefore give
\[
q_\rho(\overline A;S_{\mathrm{pair}})
\leq\frac1{|G|}\sum_{g\in G}
q_\rho(U_gAU_g^\top;S_{\mathrm{pair}})
=q_\rho(A;S_{\mathrm{pair}}).
\tag{3}
\]
Thus symmetrization does not restrict the adversary or assume the desired reduction: equation \((3)\) proves that an oracle minimizer may be sought among \(G\)-invariant covariances, and it supplies a lower-bound reduction valid for every covariance in the full class \(\mathcal C_n\).

We now derive the form of a \(G\)-invariant covariance.  A swap within pair \(j\) fixes every pair-sum vector and changes the sign of \(e_j^-\).  Commutation with all such swaps forces the cross block between the pair-sum space and \(T\) to vanish and forces all off-diagonal entries of the restriction to \(T\), in the basis \(e_1^-,\ldots,e_k^-\), to vanish.  Pair permutations then force the remaining diagonal entries on \(T\) to be equal; call their common value \(b\).  On the pair-sum space, commutation with every pair permutation forces equal diagonal entries and equal off-diagonal entries.  The all-ones direction is annihilated because every \(zz^\top\), \(z\in\mathcal Z_n\), annihilates \(\mathbf1_n\).  It follows that the restriction to the centered pair-sum space \(S_{\mathrm{pair}}\) is a scalar multiple of the identity; call that scalar \(a\).  Hence the two-eigenspace form is derived as
\[
\overline A=aP_{S_{\mathrm{pair}}}+bP_T
\quad\text{on }H_n.
\tag{4}
\]
Positive semidefiniteness gives \(a,b\geq0\).  Every member of \(\mathcal C_n\) has diagonal entries one and hence trace \(n=2k\).  Taking the trace in \((4)\) and using \((1)\) yields
\[
a(k-1)+bk=2k.
\tag{5}
\]

Put \(s=\rho^2\).  For a unit vector \(u=x+y\) according to \((1)\), put \(r=\|x\|_2^2\).  The score constraint is exactly \(r\in[s,1]\), and every such \(r\) is attainable because both subspaces in \((1)\) are nonzero.  Equations \((4)\) and \((5)\) therefore give
\[
q_\rho(\overline A;S_{\mathrm{pair}})
=\max_{r\in[s,1]}\{ar+b(1-r)\}
=\max\{a,\,as+b(1-s)\}.
\tag{6}
\]
Let \(c_n=2k/(2k-1)\).  From \((5)\),
\[
b-a=\frac{2k-(2k-1)a}{k}.
\tag{7}
\]
If \(a\geq c_n\), equation \((6)\) is at least \(a\geq c_n\).  If \(0\leq a\leq c_n\), then \(b\geq a\) by \((7)\), so \((6)\) equals its second entry.  Substitution from \((5)\) gives the affine function
\[
as+b(1-s)
=2(1-s)+a\left\{s-\frac{k-1}{k}(1-s)\right\}.
\tag{8}
\]
The coefficient of \(a\) in \((8)\) has the same sign as
\[
(2k-1)s-(k-1),
\]
and hence changes sign at
\[
s_0=\frac{k-1}{2k-1}=\frac{n-2}{2(n-1)}.
\tag{9}
\]
It follows from \((8)\) that the minimum over \(0\leq a\leq c_n\) is attained at \(a=c_n\) and equals \(c_n\) when \(s\leq s_0\), whereas it is attained at \(a=0\) and equals \(2(1-s)\) when \(s\geq s_0\).  Together with the case \(a\geq c_n\), this proves the matching lower bound
\[
q_\rho(A;S_{\mathrm{pair}})
\geq\min\{c_n,2(1-s)\}
\qquad\text{for every }A\in\mathcal C_n,
\tag{10}
\]
where equation \((3)\) is what extends the invariant calculation to the full design class.

Both endpoints used in minimizing \((8)\) are feasible.  Complete randomization has covariance \(A_{\mathrm{CRD}}=c_nP_H\), corresponding to \(a=b=c_n\), and its loss is \(c_n\).  For the within-pair design, let \(\epsilon_1,\ldots,\epsilon_k\) be independent with probability one half on each sign and set
\[
z_{2j-1}=\epsilon_j,
\qquad z_{2j}=-\epsilon_j.
\tag{11}
\]
Every assignment in its support lies in \(\mathcal Z_n\), and \(\mathbb E[z]=0\), so the law belongs to \(\mathcal D_n\).  Its covariance is block diagonal with \(k\) copies of
\[
\begin{pmatrix}1&-1\\-1&1\end{pmatrix},
\]
and hence equals \(2P_T\), corresponding to \(a=0\) and \(b=2\).  Its loss in \((6)\) is \(2(1-s)\).  These two designs attain the lower bound \((10)\), including the boundary values \(s=0\) and \(s=1\).  Therefore
\[
v^*(\rho;S_{\mathrm{pair}})
=\min\{c_n,2(1-\rho^2)\}.
\tag{12}
\]

The within-pair covariance annihilates \(S_{\mathrm{pair}}\), while every feasible covariance is positive semidefinite.  Thus
\[
\kappa(S_{\mathrm{pair}})=0.
\tag{13}
\]
By \((9)\) and \((12)\), strict improvement over complete randomization occurs exactly when \(\rho^2>s_0\).  Hence the definition of the quality threshold gives \(\rho^*(S_{\mathrm{pair}})=\sqrt{s_0}\).  The admissibility-frontier theorem identifies this threshold with \(\sqrt{\gamma(S_{\mathrm{pair}})}\); therefore
\[
\gamma(S_{\mathrm{pair}})=s_0=\frac{n-2}{2(n-1)},
\qquad
\rho^*(S_{\mathrm{pair}})=\sqrt{\frac{n-2}{2(n-1)}}.
\]
Since \(0<s_0<1\) for \(n\geq4\), this also proves \(S_{\mathrm{pair}}\in\mathfrak A_n\).  All optimizers used above are supported on balanced assignments, have mean zero, and give each unit marginal treatment probability one half.